\section{The complete off-critical metric in both phases}
\noindent
Sections \S4--\S8 obtained the \emph{singular} part of $\Tm_{\mu\nu}$ near the critical line.
Minimal paths need the whole tensor everywhere. This section supplies it: closed form in the
ordered phase, and in terms of three Watson integrals in the disordered phase. Both
expressions have been checked against the matrix solvers of \S7 and \S8 to all printed
digits.

\medskip
\noindent
\FLAG{A correction to \S8 first.} In \S8.5 the numerator of $\bar b$ was identified with
	$dF$. That is true only \emph{on} the critical line. Off it (\S\ref{sec:offbdy}),
	\begin{equation}
		\bar b=\frac{dF}{4\Gamma x^{2}}-\frac{d\beta}{4\Gamma\beta} ,
	\end{equation}
	and the second, finite term --- invisible to a leading singular analysis, since it is down by
	$x^{2}$ --- is exactly what completes the tensor. The same warning applies to the disordered
	side: $F$ and the saddle quantities agree only at $r=0$.
	
	%===============================================================
\subsection{Notation}
As in \S0--\S9. In addition, define the reduced controls
\begin{equation}
  \boxed{\ t\equiv\frac{T}{T_0}=\frac{\beta_0}{\beta},\qquad
   \eta\equiv\frac{H}{4P}=\frac{h}{4\beta P}\ }
  \label{eq:reduced}
\end{equation}
so that the critical line \eqref{eq:Fdef} of \S6 is simply the parabola
\begin{equation}
  F(t,\eta)=t+\eta^{2}-1=0 ,\qquad F<0\ \text{ordered},\ F>0\ \text{disordered}.
\end{equation}
In the ordered phase $M=\eta$ and $x^{2}=-F$. Throughout, $D\equiv a_0V=a_0I_2+2M^{2}$ with
$V=I_2+2M^{2}/a_0$ the quantity of \S3.7 and \S7.3; $D$ is the more convenient grouping here.

\subsection{Two identities that hold only on the boundary}
\label{sec:offbdy}
Both are elementary and both were silently assumed in \S6 and \S8.

\paragraph{(i) Disordered branch.} The constraint is $I_1+M^{2}=1$, so
\begin{equation}
  F=t+M^{2}-1=t-I_1 .
\end{equation}
Since $\Wc$ decreases with $r$ and $\Wc(0)=t$, we have $I_1=\Wc(r)<t$ for $r>0$, i.e.\ $F>0$
--- correct, but note $F\neq t+M^{2}-1$ is \emph{not} the same as ``$M^{2}=1-\Wc(0)$''. The
magnetisation obeys $M=h/2a_0$ with $a_0=r+2\beta P$, whereas $F$ is defined with $2\beta P$
alone. The two coincide only at $r=0$.

\paragraph{(ii) Ordered branch.} With $F=t+M^{2}-1$ and $M=h/2a_0$, $a_0=2\beta P$,
\begin{equation}
  dF=-\frac{t+2M^{2}}{\beta}\,d\beta+\frac{M}{a_0}\,dh ,
\end{equation}
whereas the numerator appearing in $\bar b$ (\S8.5) is
$(1+M^{2})s_{\Ac}+\tfrac{M}{a_0}s_G$ with $s_{\Ac}=-d\beta/\beta$, $s_G=dh$. Their difference
is
\begin{equation}
  \Big[(1+M^{2})-(t+2M^{2})\Big]\Big(-\frac{d\beta}{\beta}\Big)
  =-\frac{1-t-M^{2}}{\beta}d\beta=-\frac{x^{2}}{\beta}d\beta ,
\end{equation}
giving the boxed correction above. On the critical line $x^{2}=0$ and the two agree, which is
why \S8's singular result was right.

\subsection{Disordered phase: the general-source solution}
From \S6.5, with a general source $q_s=s_{\Ac}\delta\Ac+s_G G$, the gauge $1+2\Gamma\kappa=0$
gives $b_{\bk}=\lambda_s/4\Gamma a_{\bk}$ with
\begin{equation}
  \boxed{\ \lambda_s=-\frac{a_0(1+M^{2})s_{\Ac}+M s_G}{D}\ }
  \label{eq:lam10}
\end{equation}
What \S6 did not need, and what is required here, is the coefficient of $G$ in the solution.
Writing $\phi_s=\gamma_s G+\sum_{\bk}b_{\bk}Y_{\bk}$, we have $\gamma_s=\alpha_s+2b_0M$ with
$\alpha_s=-(s_{\Ac}+\lambda_s V)/2\Gamma M$ from \S6.5 and $b_0=\lambda_s/4\Gamma a_0$. Then
\begin{equation}
  \gamma_s=\frac{1}{2\Gamma}\Big[-\frac{s_{\Ac}}{M}-\frac{\lambda_s}{M}\Big(I_2+\frac{M^{2}}{a_0}\Big)
   +\frac{M\lambda_s}{a_0}\Big]
  =-\frac{s_{\Ac}+\lambda_s I_2}{2\Gamma M} ,
\end{equation}
the two $M\lambda_s/a_0$ terms cancelling. Substituting \eqref{eq:lam10} and using
$D-a_0I_2(1+M^{2})=M^{2}(2-a_0I_2)$,
\begin{equation}
  \boxed{\ \gamma_s=\frac{M(a_0I_2-2)s_{\Ac}+I_2\,s_G}{2\Gamma D}\ }
  \label{eq:gam10}
\end{equation}
\FLAG{Why $\gamma$ was invisible before.} Near the line $\lambda_s\to0$ like $\sqrt r$ while
$b_{\bQ}\propto\lambda_s/r$ blows up, so the singular part of $\Tm$ comes entirely from
$\lambda_s$. $\gamma_s$ stays finite and contributes only to the regular background --- which
is precisely what fast marching needs away from the line.

\subsection{Disordered phase: the bilinear form}
Conditioning the Gaussian covariance on $\delta Q_+=0$ (\S7.3) gives two compact objects.
Using $\tfrac1N\mathrm{Cov}(q_s,\delta Q)=(1+M^{2})s_{\Ac}+\tfrac{M}{a_0}s_G=-D\lambda_s/a_0$
and $\tfrac1N\mathrm{Var}(\delta Q)=V=D/a_0$:
\begin{align}
  \frac1N\avg{q_s\,G}_c&=\underbrace{Ms_{\Ac}+\frac{s_G}{2a_0}}_{\text{unconditioned}}
   \;+\;\frac{M}{a_0}\lambda_s\;\equiv\;R_s ,
  \label{eq:Rs}\\
  \frac1N\avg{q_s\,\Gc}_c&=I_2s_{\Ac}+I_3\lambda_s\;\equiv\;S_s ,
  \qquad \Gc\equiv\sum_{\bk}\frac{Y_{\bk}}{a_{\bk}} ,
  \label{eq:Ss}
\end{align}
where in each case the conditional correction $-\mathrm{Cov}(q_s,\delta Q)\mathrm{Cov}(\cdot,\delta Q)/\mathrm{Var}(\delta Q)$
collapses because the factor $-D\lambda_s/a_0$ cancels against $V=D/a_0$. Since
$\phi_t=\gamma_tG+\tfrac{\lambda_t}{4\Gamma}\Gc$,
\begin{equation}
  \boxed{\ \Tm^{\rm dis}(s,t)=\gamma_t R_s+\frac{\lambda_t}{4\Gamma}S_s\ }
  \label{eq:bilin10}
\end{equation}
Symmetry under $s\leftrightarrow t$ is guaranteed by detailed balance (\S7.4) and is a useful
algebraic check on \eqref{eq:lam10}--\eqref{eq:gam10}.

\subsection{Disordered phase: the tensor}
Insert the basis sources $e_\beta:(s_{\Ac},s_G)=(-1/\beta,0)$ and $e_h:(0,1)$, for which
\begin{equation}
  \lambda_\beta=\frac{a_0(1+M^{2})}{\beta D},\quad
  \gamma_\beta=-\frac{M(a_0I_2-2)}{2\Gamma\beta D},\qquad
  \lambda_h=-\frac{M}{D},\quad
  \gamma_h=\frac{I_2}{2\Gamma D} .
\end{equation}
Expanding \eqref{eq:bilin10} and collecting powers of $I_2,I_3$:
\begin{equation}
  \boxed{
  \begin{aligned}
  \Tm^{\rm dis}_{\beta\beta}&=\frac{a_0^{2}\big[(1+M^{2})^{2}I_3-(1-M^{2})I_2^{2}\big]
    -8a_0M^{2}I_2+4M^{2}(1-M^{2})}{4\Gamma\beta^{2}D^{2}},\\[2pt]
  \Tm^{\rm dis}_{\beta h}&=-\frac{M\big[a_0I_2^{2}-2I_2+a_0(1+M^{2})I_3\big]}{4\Gamma\beta D^{2}},\\[2pt]
  \Tm^{\rm dis}_{hh}&=\frac{I_2^{2}+M^{2}I_3}{4\Gamma D^{2}} .
  \end{aligned}}
  \label{eq:Tdis}
\end{equation}
The determinant factorises,
\begin{equation}
  \det\Tm^{\rm dis}=\frac{(1-M^{2})I_3-I_2^{2}}{16\Gamma^{2}\beta^{2}D^{2}}\;>\;0 ,
  \label{eq:detdis}
\end{equation}
positivity following from Cauchy--Schwarz, $I_2^{2}\le I_1I_3=(1-M^{2})I_3$ (\S4.7). This is a
stronger statement than anything in \S6: the tensor is \emph{full rank everywhere} in the open
disordered phase, not merely near the line.

\paragraph{Check.} At $h=0$, $M=0$, $D=a_0I_2$, and \eqref{eq:Tdis} collapses to
\begin{equation}
  \Tm^{\rm dis}_{\beta\beta}=\frac{1}{4\Gamma\beta^{2}}\Big(\frac{I_3}{I_2^{2}}-1\Big),
  \qquad \Tm^{\rm dis}_{\beta h}=0,\qquad \Tm^{\rm dis}_{hh}=\frac{1}{4\Gamma a_0^{2}} ,
\end{equation}
which is exactly \S4.7 and \S4.9. \checkmark

\subsection{Ordered phase: the tensor}
Here \S8 already supplies everything except the expansion. From \S8.5--\S8.6,
\begin{equation}
  \alpha_s=\frac{1}{2\Gamma}\Big[Ms_{\Ac}+\frac{s_G}{a_0}\Big],
  \qquad
  \bar b_s=\frac{(1+M^{2})s_{\Ac}+(M/a_0)s_G}{4\Gamma x^{2}} ,
\end{equation}
\begin{equation}
  \Tm^{\rm ord}(s,t)=\alpha_t\Big[Ms_{\Ac}+\frac{s_G}{2a_0}\Big]
   +\bar b_t\Big[(t+2M^{2})s_{\Ac}+\frac{M}{a_0}s_G\Big] .
  \label{eq:bilinord}
\end{equation}
Inserting the basis sources gives, after using $x^{2}=1-t-M^{2}$ three times,
\begin{equation}
  \boxed{\ \Tm^{\rm ord}_{(\beta,h)}=\frac{1}{4\Gamma x^{2}}
  \begin{pmatrix}
   \dfrac{t+(4-t)M^{2}}{\beta^{2}} & -\dfrac{M(2-t)}{a_0\beta}\\[10pt]
   -\dfrac{M(2-t)}{a_0\beta} & \dfrac{1-t}{a_0^{2}}
  \end{pmatrix}\ }
  \label{eq:Tord}
\end{equation}
The three simplifications are worth recording, since they are what makes the result compact:
\begin{align}
  2M^{2}x^{2}+(1+M^{2})(t+2M^{2})&=t+(4-t)M^{2},\\
  2x^{2}+t+2M^{2}&=2-t,\\
  x^{2}+M^{2}&=1-t .
\end{align}
The determinant is
\begin{equation}
  \det\Tm^{\rm ord}=\frac{t}{16\Gamma^{2}\beta^{2}a_0^{2}x^{2}}>0 ,
\end{equation}
using $\big(t+(4-t)M^{2}\big)(1-t)-M^{2}(2-t)^{2}=t\,x^{2}$.

\FLAG{No null directions, in either phase.} \eqref{eq:detdis} and the above are both strictly
positive. The rank-one degeneracy that drives the optimal-path classification in [BR,\S e] is
\emph{not} inherited by the spherical model --- consistent with the static statement in
\S3.10, and now established for the full dynamical tensor.

\subsection{Reduced coordinates, and an exact decomposition}
The Jacobian from \eqref{eq:reduced} is
\begin{equation}
  \frac{\partial(\beta,h)}{\partial(t,\eta)}=
  \begin{pmatrix}-\beta/t & 0\\ -h/t & 4\beta P\end{pmatrix},
\end{equation}
and pulling \eqref{eq:Tord} back with $\Tm^{(y)}=J^{\top}\Tm^{(\lambda)}J$, using
$h/a_0=2\eta$ and $4\beta P=2a_0$, gives
\begin{equation}
  \boxed{\ \Tm^{\rm ord}_{(t,\eta)}=\frac{1}{4\Gamma x^{2}}
  \begin{pmatrix}\dfrac{1-\eta^{2}}{t} & 2\eta\\[6pt] 2\eta & 4(1-t)\end{pmatrix}\ }
  \label{eq:Tred}
\end{equation}
and --- the most useful form for computation --- the \emph{exact} decomposition
\begin{equation}
  \boxed{\ d\ell_T^{2}\Big|_{\rm ord}
   =\frac{(dF)^{2}}{4\Gamma x^{2}}+\frac{(dt)^{2}}{4\Gamma t}+\frac{(d\eta)^{2}}{\Gamma}\ }
  \label{eq:decomp}
\end{equation}
valid throughout the ordered phase, not just near the boundary. To verify, use
$dF=dt+2\eta\,d\eta$ and $t+x^{2}=1-\eta^{2}$, $x^{2}+\eta^{2}=1-t$.

\paragraph{Reading.} \eqref{eq:decomp} separates the geometry into three pieces: a
\emph{singular normal} term $\propto(dF)^{2}/x^{2}$, which is the rank-one structure of \S8
and diverges as $|F|^{-1}$; a \emph{thermal background} $(dt)^{2}/4\Gamma t$, which is finite
everywhere except $T\to0$; and a \emph{field background} $(d\eta)^{2}/\Gamma$, which is flat.
The first was all that \S8 saw. The second and third are the regular part, and they are what
decide the ``around versus through'' competition away from the line.

Note also that $(dt)^{2}/4\Gamma t=\big(d\sqrt t\big)^{2}/\Gamma$, so in coordinates
$(2\sqrt t,\,2\eta)$ the two background terms together are \emph{flat Euclidean} --- the same
combination that appeared, spuriously, in the very first draft of these notes. It is real
here, but it is only part of the metric; the singular normal term is not flat. Remarkably,
adding that term back does not spoil matters but produces a different constant-curvature
geometry, as the next subsection shows.

\subsection{The ordered phase is a spherical patch --- and the sphere is the constraint}
\label{sec:sphere}
There is a three-line route to this, and it is more revealing than the coordinate
manipulation that follows. Take the exact decomposition \eqref{eq:decomp} and define
\begin{equation}
  \boxed{\ X\equiv\sqrt{t}=\sqrt{\Wc},\qquad Y\equiv\eta=M,\qquad Z\equiv x\ }
  \label{eq:XYZ}
\end{equation}
i.e.\ the square roots of the three terms in the spherical constraint of \S2.3. That
constraint, $\Wc+M^{2}+x^{2}=1$, is then literally
\begin{equation}
  X^{2}+Y^{2}+Z^{2}=1 .
\end{equation}
Now $F=-x^{2}$ gives $dF=-2x\,dx$, so $(dF)^{2}/4\Gamma x^{2}=dZ^{2}/\Gamma$; and
$t=X^{2}$ gives $dt^{2}/4\Gamma t=dX^{2}/\Gamma$; and $d\eta^{2}/\Gamma=dY^{2}/\Gamma$.
Hence
\begin{equation}
  \boxed{\ d\ell_T^{2}\Big|_{\rm ord}=\frac{1}{\Gamma}\Big[dX^{2}+dY^{2}+dZ^{2}\Big]
   \quad\text{restricted to}\quad X^{2}+Y^{2}+Z^{2}=1\ }
  \label{eq:sphereXYZ}
\end{equation}
the round metric induced on a sphere of radius $R=1/\sqrt\Gamma$.

\FLAG{This is not a coordinate accident.} The three Cartesian coordinates are the thermal
occupation, the uniform magnetisation and the staggered order parameter, and the sphere they
lie on \emph{is the spherical constraint that defines the model}. The thermodynamic geometry
of the ordered phase is the round metric on the constraint sphere itself. That the
dissipation metric --- built from the dynamics, the projected generator and a resolvent ---
should land exactly on the model's own configuration-space constraint is a statement worth
making explicitly in the manuscript.

\paragraph{Equivalent route in angles.} For completeness, and because the angular
coordinates are what a solver wants, set
\begin{equation}
  w\equiv2\sqrt t,\qquad v\equiv2\eta,
  \qquad\text{so}\qquad
  \frac{(dt)^{2}}{4\Gamma t}+\frac{(d\eta)^{2}}{\Gamma}=\frac{dw^{2}+dv^{2}}{4\Gamma},
\end{equation}
and introduce polar coordinates on the $(w,v)$ plane,
\begin{equation}
  w=\varrho\cos\psi,\qquad v=\varrho\sin\psi,\qquad \varrho^{2}=w^{2}+v^{2} .
\end{equation}
Since $F=t+\eta^{2}-1=\tfrac14(\varrho^{2}-4)$ we have $x^{2}=-F=\tfrac14(4-\varrho^{2})$ and
$dF=\tfrac12\varrho\,d\varrho$, so the singular normal term becomes
\begin{equation}
  \frac{(dF)^{2}}{4\Gamma x^{2}}=\frac{\varrho^{2}\,d\varrho^{2}}{4\Gamma\,(4-\varrho^{2})} .
\end{equation}
Adding the background, the cross terms cancel and \eqref{eq:decomp} becomes \emph{diagonal}:
\begin{equation}
  d\ell_T^{2}\Big|_{\rm ord}
   =\frac{1}{4\Gamma}\left[\Big(1+\frac{\varrho^{2}}{4-\varrho^{2}}\Big)d\varrho^{2}
    +\varrho^{2}d\psi^{2}\right]
   =\frac{1}{4\Gamma}\left[\frac{4\,d\varrho^{2}}{4-\varrho^{2}}+\varrho^{2}d\psi^{2}\right].
  \label{eq:polar}
\end{equation}
Finally substitute $\varrho=2\sin\chi$, for which $4-\varrho^{2}=4\cos^{2}\chi$ and
$d\varrho=2\cos\chi\,d\chi$:
\begin{equation}
  \boxed{\ d\ell_T^{2}\Big|_{\rm ord}=\frac{1}{\Gamma}\Big[d\chi^{2}+\sin^{2}\!\chi\;d\psi^{2}\Big]\ }
  \label{eq:sphere}
\end{equation}
\textbf{Each selected ordered pure-state branch is isometric to a quarter of a round
two-sphere of radius $R=1/\sqrt\Gamma$}, with constant \emph{Gaussian} (sectional) curvature
\begin{equation}
  K=\frac{1}{R^{2}}=\Gamma
  \qquad\text{(the Ricci scalar is }\mathcal R=2K=2\Gamma\text{)} .
\end{equation}
Verified numerically: transforming \eqref{eq:Tred} to $(\chi,\psi)$ by the Jacobian of
$(t,\eta)\leftrightarrow(\chi,\psi)$ reproduces $\Gamma^{-1}\mathrm{diag}(1,\sin^{2}\chi)$ to
eight digits at every test point.

\FLAG{A patch, not the whole sphere.} The selected branch has $X=\sqrt t>0$ and $Z=x>0$, so
the domain is the intersection of two hemispheres --- a quarter-sphere. The second
symmetry-broken branch $x<0$ is the reflected quarter, and if both state sheets are retained
their union is the hemisphere $X>0$, with the critical line $x=0$ the great circle separating
them. Saying flatly ``the ordered phase is a round two-sphere'' overstates it.

\paragraph{The domain.} Since $t>0$ requires $w>0$ and the ordered phase requires
$\varrho<2$,
\begin{equation}
  \chi\in[0,\tfrac\pi2),\qquad \psi\in(-\tfrac\pi2,\tfrac\pi2),
\end{equation}
a lune covering half the northern hemisphere. The boundaries all lie in the smooth metric
completion of the ordered phase (they are not in the open phase itself), and in that
completion they are geodesics:
\begin{center}
\begin{tabular}{lll}
\hline
locus & sphere image & is a geodesic?\\
\hline
critical line $F=0$ & equator $\chi=\pi/2$ & \textbf{yes} (great circle)\\
$T=0$ & meridians $\psi=\pm\pi/2$ & yes\\
$T=0,H=0$ & pole $\chi=0$ & (a point)\\
\hline
\end{tabular}
\end{center}

\paragraph{Consequences.} Three, and all matter for \S11.
\begin{enumerate}
  \item \textbf{Ordered-phase geodesics are great-circle arcs}, and distances are
        $\ell=\Theta/\sqrt\Gamma$ with $\Theta$ the great-circle angle. No numerical solver
        is needed inside the ordered phase at all.
  \item \textbf{The critical line is itself a geodesic.} A minimal path that reaches the
        boundary may run \emph{along} it at locally minimal cost. This is a third candidate
        topology, alongside ``around the tip'' and ``straight through'', and it was not
        visible before \eqref{eq:sphere}: the singularity of the metric is purely normal
        (\S8.6), so the tangential direction along $F=0$ is perfectly regular.
  \item \textbf{An exact benchmark for the numerics.} Great-circle distances give a
        non-trivial closed-form answer against which any eikonal or graph solver can be
        validated before it is trusted in the disordered phase.
\end{enumerate}

\subsection{The [BR] convention is hyperbolic, not spherical}
\label{sec:hyperbolic}
\FLAG{Correction.} An earlier version of this section asserted that the convention factor
``destroys constant curvature''. That is wrong for the choice actually made in [BR].

Pass to the Mercator coordinate $d\mu=d\chi/\sin\chi$, i.e.\ $\mu=\ln\tan(\chi/2)$ and
$\sin\chi=\operatorname{sech}\mu$, in which
$d\chi^{2}+\sin^{2}\!\chi\,d\psi^{2}=\sin^{2}\!\chi\,(d\mu^{2}+d\psi^{2})$. With
$t=\sin^{2}\!\chi\cos^{2}\!\psi$ and $\beta=\beta_0/t$, the [BR] metric $g=\beta\Tm$ is
\begin{equation}
  d\ell_g^{2}
  =\frac{\beta_0}{\sin^{2}\!\chi\cos^{2}\!\psi}\cdot\frac{\sin^{2}\!\chi}{\Gamma}
   \big(d\mu^{2}+d\psi^{2}\big)
  =\boxed{\ \frac{\beta_0}{\Gamma}\,\frac{d\mu^{2}+d\psi^{2}}{\cos^{2}\!\psi}\ }
  \label{eq:poincare}
\end{equation}
\textbf{The factor $\sin^{2}\!\chi$ cancels exactly}, leaving the Poincar\'e metric on the
strip $|\psi|<\pi/2$. Its Gaussian curvature is constant and \emph{negative},
\begin{equation}
  \boxed{\ K_g=-\frac{\Gamma}{\beta_0}\ },\qquad \mathcal R_g=-\frac{2\Gamma}{\beta_0},
\end{equation}
(or $K_g=-\Gamma$ if the dimensionless factor $\beta/\beta_0$ is used). So
\begin{equation}
  \boxed{\ \Tm^{\rm ord}\ \text{is spherical};\qquad
   g^{\rm BR}=\beta\,\Tm^{\rm ord}\ \text{is hyperbolic.}\ }
\end{equation}
Both are exactly solvable. \emph{No numerical solver is needed inside the ordered phase for
either convention} --- one simply uses hyperbolic geodesics instead of great circles.

\paragraph{The boundary survives the change.} The critical line $\chi=\pi/2$ is $\mu=0$, and
since \eqref{eq:poincare} is independent of $\mu$, the reflection $\mu\to-\mu$ is an isometry
whose fixed-point set is $\mu=0$. A fixed-point set of an isometry is totally geodesic, so
the critical line is a geodesic of the [BR] metric too. The third path topology of
\S\ref{sec:sphere} therefore survives the convention change, even though the great-circle
description does not.

\FLAG{Only two of the three conventions are constant-curvature.} Computing
$K=-e^{-2\phi}\Delta\phi$ numerically for $d\ell^{2}=e^{2\phi}(d\mu^{2}+d\psi^{2})$:
\begin{center}
\begin{tabular}{lcccc}
\hline
$(\mu,\psi)$ & $(-0.4,0.2)$ & $(-1.0,-0.6)$ & $(-0.2,0.9)$ & $(-2.0,0.1)$\\
\hline
$K$ for $\Tm$            & $+1.0000$ & $+1.0000$ & $+1.0000$ & $+1.0000$\\
$K$ for $\beta\Tm$ (BR)  & $-1.0000$ & $-1.0000$ & $-1.0000$ & $-1.0000$\\
$K$ for $\Tm/\beta$      & $3.914$   & $19.21$   & $12.64$   & $232.99$\\
\hline
\end{tabular}
\end{center}
($\Gamma=\beta_0=1$.) The excess-work pullback $\Tm/\beta$ --- the convention that the tensor
transformation of \S4.11 actually favours --- gives \emph{no} constant curvature. This is a
genuine tension: the convention with the clean geometry is not the one the Sivak--Crooks
transformation produces. Elegance is not evidence, so \S4.11 remains open; but the
geometric signature is sharp enough that it should be recorded alongside the argument
there.

In all three cases the metric is conformally flat in $(\mu,\psi)$, so the ordered-phase
problem is \emph{isotropic} in Mercator coordinates and an ordinary scalar-speed fast marcher
suffices if one is used at all.

\subsection{Numerical assembly}
For a grid point $(t,\eta)$:
\begin{enumerate}
  \item $F=t+\eta^{2}-1$.
  \item If $F<0$: use \eqref{eq:Tred} directly with $x^{2}=-F$, $M=\eta$ --- no root solve.
        Better still, skip the grid entirely: the ordered phase is exactly solvable in
        \emph{both} constant-curvature conventions --- great circles for $\Tm$
        (\S\ref{sec:sphere}), hyperbolic geodesics on the Poincar\'e strip for $g^{\rm BR}$
        (\S\ref{sec:hyperbolic}).
  \item If $F>0$: solve the saddle equation $I_1(r,\beta)+M^{2}=1$ for $r$ at
        $\beta=\beta_0/t$, $M=h/2a_0$, then evaluate $I_2,I_3$ and use \eqref{eq:Tdis},
        pulling back with the Jacobian above.
  \item Apply the scalar convention factor of \S4.11 --- and see the warning there, which is
        \emph{not} resolved by this section.
\end{enumerate}
The Watson integrals depend on $(r,\beta)$ only through $s=r/\beta\Jc$, so
$I_n=(\beta\Jc)^{-n}K_n(s)$ with $K_n$ universal; tabulating $K_{1,2,3}(s)$ once and
interpolating turns each grid point into a scalar root solve.

\subsection{The critical interface}
The raw tensor is infinite on $F=0$, so a grid value placed exactly there creates a spurious
numerical barrier even though the distance is finite. Two clean treatments:
\begin{itemize}
  \item \textbf{Signed square-root coordinate} $\rho=\mathrm{sgn}(F)\sqrt{|F|}$, so
        $F=\rho|\rho|$ and $dF=2|\rho|d\rho$. Then
        \begin{equation}
          \frac{(dF)^{2}}{8\Gamma F}=\frac{(d\rho)^{2}}{2\Gamma}\ (\rho>0),
          \qquad
          \frac{(dF)^{2}}{4\Gamma|F|}=\frac{(d\rho)^{2}}{\Gamma}\ (\rho<0),
        \end{equation}
        both finite, with the ratio $2$ of \S8.6 now visible as a jump in a \emph{regular}
        coefficient.
  \item \textbf{Analytic edge cost} for an edge crossing nearly normally between $F_+>0$ and
        $F_-<0$:
        $\Delta\ell\simeq\sqrt{F_+/2\Gamma}+\sqrt{|F_-|/\Gamma}$, adding the regular part by
        midpoint quadrature.
\end{itemize}
\FLAG{And note what the interface \emph{is}.} By \S\ref{sec:sphere} and
\S\ref{sec:hyperbolic} the critical line is a geodesic of the ordered-phase metric in both
constant-curvature conventions, so it is not merely a place where the grid needs care:
it is a locus along which a minimal path may legitimately travel. Any regularisation must
therefore represent the tangential direction accurately, not just avoid the normal
divergence.

Since $\Tm_{t\eta}\neq0$ in general, the eikonal problem
$\nabla D^{\top}\Tm^{-1}\nabla D=1$ is anisotropic: a scalar-speed fast marcher is not
adequate; use an ordered-upwind or anisotropic stencil taking the full matrix.

\subsection{Verification}
Both tensors were checked against the independent matrix solvers of \S7 (disordered:
generator \eqref{eq:Nop} plus conditioned covariance, augmented solve) and \S8 (ordered:
generator plus diagonal covariance, direct solve), on a $6^{3}$ mode grid:
\begin{center}
\begin{tabular}{llccc}
\hline
phase & point & $\Tm_{\beta\beta}$ & $\Tm_{\beta h}$ & $\Tm_{hh}$\\
\hline
dis & $\beta{=}0.30,h{=}0.3$ & $0.801357$ & $-0.043512$ & $0.059339$\\
dis & $\beta{=}0.40,h{=}2.0$ & $1.426959$ & $-0.169025$ & $0.037516$\\
dis & $\beta{=}0.30,h{=}4.0$ & $0.389440$ & $-0.063877$ & $0.012418$\\
ord & $t{=}0.55,M{=}0.40$    & $1.93878$  & $-0.17007$  & $0.02199$\\
ord & $t{=}0.20,M{=}0.80$    & $8.39286$  & $-0.76531$  & $0.07086$\\
\hline
\end{tabular}
\end{center}
Closed forms and solvers agree to every digit shown, and \eqref{eq:detdis} reproduces the
directly computed determinant.

\subsection{What this changes, and what it does not}
\paragraph{Changes.} The metric is now available everywhere in both phases, so \S11 (minimal
paths) is unblocked. Three consequences for [BR]:
\begin{enumerate}
  \item \textbf{No null directions} anywhere, so the constant-magnetisation zero-length paths
        of the mean-field example have no analogue and the four-class path taxonomy of
        [BR,\S e] must be re-derived.
  \item \textbf{The whole domain must be solved.} There is no collapse of the disordered
        phase to a one-dimensional quotient.
  \item \textbf{Around versus through is a genuine quantitative competition}, not a
        consequence of a rank-one artefact.
  \item \textbf{A third topology exists.} Because the critical line is a geodesic
        (\S\ref{sec:sphere}), ``enter the ordered phase and run along the boundary'' is a
        candidate distinct from both ``around the tip'' and ``straight through''. The
        classification in [BR,\S e] has no analogue of it.
  \item \textbf{Half the problem is analytic, in either convention.} The ordered phase needs
        no numerics: geodesics are great-circle arcs for $\Tm$ and hyperbolic geodesics for
        $g^{\rm BR}$, so the crossing problem reduces to matching two boundary points on
        $F=0$ to the disordered-phase solution.
  \item \textbf{A geometric signature on the convention question.} $\Tm$ and $\beta\Tm$ both
        give constant curvature ($+\Gamma$ and $-\Gamma/\beta_0$); $\Tm/\beta$ does not
        (\S\ref{sec:hyperbolic}). This does not settle \S4.11 --- the tensor transformation
        there favours $\Tm/\beta$ --- but it is a sharp structural fact that the manuscript
        should weigh.
\end{enumerate}

\paragraph{Does not change.} The three caveats of \S4.11, \S9 and \S12 stand:
the mobility choice is not determined by the equilibrium model; the $Q_-$ correction affects
the finite-size window (\S9) but not the continuum tensor above; and the objective functional
is still unfixed, which matters here more than before --- a non-constant scalar factor is a
conformal change and \emph{changes geodesics}, so it must be settled before any path is
reported.

\TODO{\S11: anisotropic fast marching on \eqref{eq:Tred}/\eqref{eq:Tdis}.}

	